\subsubsection{Conclusioni sull'implementazione dei Covert Channel}
Questo lavoro di tesi ha affrontato lo sviluppo e l'implementazione di un 
covert channel basato sul protocollo ICMP per l'estrazione di dati. 
L'obiettivo principale era dimostrare come il protocollo ICMP, fondamentale 
per il funzionamento delle reti ma spesso sottovalutato dal punto di vista 
della sicurezza, possa essere sfruttato per creare canali di comunicazione 
nascosti efficaci.
\vspace{3ex} \newline 
Durante lo sviluppo sono state implementate diverse tipologie di covert channel:
\begin{itemize}
    \item Storage Covert Channel implementato sfruttando i vari campi dei 
    messaggi ICMP. 
    \item Timing Covert Channel codificando le informazioni negli intervalli 
    temporali tra i pacchetti. Per poter inviare un maggior numero di informazioni 
    (possibilmente senza renderlo maggiormente sensibile agli errori), si sono 
    sviluppate varie varianti per il calcolo degli intervalli di tempo. 
    \item Behavioural Covert Channel che codifica i dati attraveso l'invio di 
    una specifica tipologia di messagio e codice ICMP. 
    \item Un implementazione ibrida dei Covert Channel precedentemente fatti. 
\end{itemize}
\vspace{4ex}  
Alcune difficoltà sono state riscontrate durante l'implementazione della 
comunicazione fra le varie entità. Non potendo affidarsi a TCP, ma solo a 
messaggi ICMP, si sono dovuti gestire metodi per stabilire e rilasciare la 
connessione stabilita. Oltre all'accettarsi che la comunicazione sia 
effettivamente stabilita. 
\vspace{4ex}  \newline 
Ulteriori difficoltà tecniche sono stati l'invio di messaggi non conformi 
agli standard RFC. Si sono dovute trovare alternative per poter inserire 
dati dove da linee guida non è possibile e gestire la presenza di tipologie 
di messaggi ICMP deprecati. Inoltre siccome l'attacco deve poter essere 
eseguibile anche tramite IPv6; si sono dovuti implementare metodi per la 
ricerca dell'interfaccia da utilizzare per l'invio dei dati. 

\subsubsection{Conclusioni sui test effettuati}
Nella trasmisisone che invia una singola volta i dati, RITA non ha rilevato 
la presenza di un covert Channel. Ma questo caso non è applicabile alla 
realtà; in cui l'attaccante e la vittima si possono scambiare più volte 
i dati. Si è quindi proceduto a testare molteplici invvi dei dati. 
\vspace{2ex} \newline 
Le conclusioni sono che l'utilizzo di un delay fra l'invio di un blocco di 
dati ed il successivo permette di abbassare la visibilità del Covert Channel. 
Anche la falsificazione dei mittenti permette questa cosa, a patto che gli 
indirizzi IP utilizzabili non vengano utilizzati in sequenza per ogni pacchetto 
inviato ma a blocchi. Così da poter permettere agli IP precedentemente utilizzati 
di abbassare la loro visibilità. 
\vspace{1ex} \newline
In sintesi la rilevabilità di un Covert channel dipende da: 
\begin{itemize}
    \item Frequenza di trasmissione: 
    Trasmissioni sporadiche sono più difficili da rilevare rispetto a 
    comunicazioni con pattern regolari. 
    \item Distribuzione delle sorgenti: 
    L'utilizzo di mittenti multipli (reali o falsificati intelligentemente), 
    permette di ridurre significativamente la rilevabilità. 
    \item Intervalli irregolari e sufficientemente lunghi tra le 
    trasmissioni riducono la probabilità di essere identificati come 
    beaconing
\end{itemize}

\subsubsection{Conclusioni sull'utilizzo di ICMP}
Sebbene sia un protocollo funzionale al monitoraggio di rete, i suoi messaggi 
devono essere limitati a quelli essenziali. In particolare le principali 
tipologie di messaggi ICMP che servono sono la tipologia Echo Request/Reply 
(usata in strumenti come ping), la tipologia Destination IUnreachable e la 
tipologia Time Exceeded. Il resto delle tipologie sono state deprecate siccome 
sostituite da strumenti diversi. 
\vspace{1ex} \newline 
Inoltre fra quelli autorizzati bisogna monitorare il loro utilizzo ed il 
loro contenuto. Un messaggio Echo Request può contenere dati nel campo payload 
e messagi come Destination Unreachable o Time Exceeded possono generare 
dei pattern temporali. 
\vspace{2ex} \newline 
Ciò richiede un ispezione dei pacchetti presenti nel flusso, un analisi 
statistica per la ricerca di pattern e la limitazione (nella rete) di 
messaggi ICMP.  


